%%% FIELD proof-support-equality-reduction
Assume that \(D_{jh}=D_{jK}\) for every \((j,h)\in\mathcal E\).  By the definition of the two support-restricted feasible sets, for each \((j,h)\in\mathcal E\),
\[
 \Delta(D_{jh})=\Delta(D_{jK}).
\]
The set \(\mathcal E\) is finite because \(\mathcal J\subseteq\{1,\ldots,N\}\) and \(h\in\{0,\ldots,K\}\).  Taking the finite Cartesian product of the preceding equalities therefore gives
\[
 \mathcal F_{\mathrm{nested}}
 =\prod_{(j,h)\in\mathcal E}\Delta(D_{jh})
 =\prod_{(j,h)\in\mathcal E}\Delta(D_{jK})
 =\mathcal F_{\mathrm{deep}}.
 \tag{1}
\]

For \(s\in\{\mathrm{nested},\mathrm{deep}\}\), let \(\mathcal B_s\) denote the set of all pairs \((w,\ell)\) satisfying the constraints in \texttt{def:exact-width-frontier}.  In that definition, the support label \(s\) enters only through the constraint \(w\in\mathcal F_s\).  The model class, the deterministic-outcome restriction on \(w\) and \(\ell\), the estimator \(\widehat\tau(w)\), the causal target \(\tau\), and the uniform joint-coverage constraint are identical under the two labels.  Hence (1) implies
\[
 \mathcal B_{\mathrm{nested}}=\mathcal B_{\mathrm{deep}}.
 \tag{2}
\]
Both exact-frontier problems apply the same objective \((w,\ell)\mapsto\max_{e\in\mathcal E}\ell_e\) to these sets.  Thus \texttt{def:exact-width-frontier} and (2) yield
\[
 \begin{aligned}
 W_{\mathrm{nested}}
 &=\inf_{(w,\ell)\in\mathcal B_{\mathrm{nested}}}
     \max_{e\in\mathcal E}\ell_e\\
 &=\inf_{(w,\ell)\in\mathcal B_{\mathrm{deep}}}
     \max_{e\in\mathcal E}\ell_e
 =W_{\mathrm{deep}}.
 \end{aligned}
 \tag{3}
\]
This equality also holds under the extended-value convention if the common admissible set is empty.

Similarly, for each support label \(s\), let \(\mathcal S_s\) be the set of tuples \((w,u,v,g)\) satisfying the constraints in \texttt{def:bonferroni-socp}.  Again the only label-dependent constraint is \(w\in\mathcal F_s\); the predictor-balance norms, Gaussian-standard-error norms, scalar envelope inequalities, fixed constants, and coordinate set are otherwise the same.  Equation (1) consequently gives
\[
 \mathcal S_{\mathrm{nested}}=\mathcal S_{\mathrm{deep}}.
 \tag{4}
\]
Since the two programs minimize the same scalar objective \(g\), \texttt{def:bonferroni-socp} and (4) imply
\[
 \begin{aligned}
 \Gamma_{\mathrm{nested}}
 &=\min_{(w,u,v,g)\in\mathcal S_{\mathrm{nested}}}g\\
 &=\min_{(w,u,v,g)\in\mathcal S_{\mathrm{deep}}}g
 =\Gamma_{\mathrm{deep}}.
 \end{aligned}
 \tag{5}
\]
Equations (3) and (5) prove the two asserted equalities.  The standing condition \(\mathcal J\ne\varnothing\) is used only because both frontier constructions are stated on such inputs.  No covariance rank or eigenvalue bound, positivity beyond the definitions' domains, overlap condition, panel-boundary convention, Gaussian distributional property, or limiting argument is used.
